\chapter{The Descartes Universe}
\label{chap:descartes}

% ── The Descartes condition ──────────────────────────────────────────────────

\begin{definition}[Descartes condition]
\label{def:descartes-condition}
A quantum observable universe $(\mathcal{C}, A, \omega)$ over~$R$
(Definition~\ref{def:quantum-observable-universe},
Theorem~\ref{thm:qu-to-qou}) satisfies the \emph{Descartes condition} if the
categorical Euler characteristic
\[
  \chi(A) \;=\; \mathrm{tr}(\mathrm{id}_A) \;=\; -6 \;\in\; R.
\]
This requires $6 \in R^{\times}$ (i.e.\ $\mathrm{char}(R) \nmid 6$).  The
class $[A] \in K_0(\mathcal{C})$ has rank $-6$, which by
Theorem~\ref{thm:qu-to-qou}(3--4) forces $\kappa_n \neq 0$ for all $n$ and
produces $6$ branch points in the geometric descent (with total ramification
$R_{\mathrm{tot}} = 2d + 6$).  A quantum observable universe satisfying the
Descartes condition is called a \emph{Descartes universe}.
\end{definition}

% ── Implications on the three examples ───────────────────────────────────────

\begin{remark}[The Descartes condition on examples]
\label{rem:descartes-calabi-yau}
\label{rem:descartes-ale}
\label{rem:descartes-affine}
The constraint $\chi(E) = -6$ specialises to each archimedean example:
\begin{itemize}
  \item \emph{Calabi--Yau} (Example~\ref{ex:calabi-yau}): constrains
        the alternating sum of Betti numbers; only manifolds admitting
        a principal $\mathrm{SU}(2)$-bundle with this Euler characteristic
        realise the Descartes universe.
  \item \emph{ALE links} (Example~\ref{ex:ale}): selects finite subgroups
        $\Gamma \subset \mathrm{SU}(2)$ via arithmetic constraints on
        the McKay quiver.
  \item \emph{Affine space} $M = \mathbb{C}^2$: forces non-trivial
        bundle topology; the $q$-harmonic zero-modes of~$\Delta_q$
        reproduce the CMB thermal spectrum in the $q \to 1$ regime.
\end{itemize}
\end{remark}

% ── Descent to one complex dimension ─────────────────────────────────────────

\begin{remark}[Descent to one complex dimension]
\label{rem:descent-one-dim}
By the descent theorem (Proposition~\ref{prop:jet-prolongation}), a
two-complex-dimensional Descartes universe $(M, G, P, E)$ with $M = \mathbb{C}^2$
descends to a one-complex-dimensional variety $M_0$.  The stability of the
Euler characteristic under descent and lifting
(Proposition~\ref{prop:euler-char-stability}) ensures that the descended
universe retains $\chi(E) = -6$.  We may therefore study the Descartes
universe in the setting of one-complex-dimensional varieties.
The descent may continue further to dimension zero
(Remark~\ref{rem:descent-reduction-zero}), where the entire quantum
observable universe is captured by fibre data at a single point.
\end{remark}

% ── Genus and Euler characteristic ───────────────────────────────────────────

\begin{proposition}[Genus of the one-complex-dimensional descent]
\label{prop:genus-maximal-descent}
Let $(M, G, P, E)$ be a Descartes universe that descends to a compact
one-complex-dimensional variety $M_0$.  Then $M_0$ is a Riemann surface of
genus $g = 4$.
\end{proposition}

\begin{proof}
The Descartes condition gives $\mathrm{tr}(\mathrm{id}_A) = -6$.  Under the fibre functor~$\omega$ of the archimedean realisation, this maps to the sheaf Euler characteristic $\chi(V) = -6$, where $V = P \times_{\mathrm{SU}(2)} K^2$ is the associated rank-$2$ vector bundle over~$M_0$.  Since $G = \mathrm{SU}(2)$, the determinant bundle $\det(V)$ is trivial, so $\deg(V) = 0$.  By Riemann--Roch on the compact Riemann surface~$M_0$: $\chi(V) = \deg(V) + \mathrm{rank}(V)(1-g) = 2(1-g)$.  Setting $2(1-g) = -6$ gives $g = 4$.
\end{proof}

% ── Leptons and chirality ────────────────────────────────────────────────────

\begin{remark}[Electrons, muons and the Reeb decomposition]
\label{rem:descartes-leptons}
The Reeb decomposition (Definition~\ref{def:reeb-decomposition}) applied to
a Descartes universe yields two classes of fermionic excitations
(Definition~\ref{def:fermion}).  The degree of the vertical component
$\alpha_V$ distinguishes them: forms with $\deg\alpha_V = 1$ correspond to
electrons, while forms with $\deg\alpha_V = 2$ correspond to muons.  More
generally, each additional unit of vertical degree is identified with a
successive lepton generation.  The constraint $\chi(E) = -6$ bounds the
number of such generations, since the Hodge decomposition of $E$ admits only
finitely many independent vertical harmonic forms.  This
generation-dependent structure leads to testable predictions, including a
$q$-correction to the proton charge radius that differs between muonic and
electronic hydrogen (Section~\ref{ssec:proton-radius}).
\end{remark}

\begin{remark}[Chirality as Reeb orientation]
\label{rem:descartes-chirality}
The Reeb vector field $\xi$ on the CR manifold $M$ carries a natural
orientation: the sign of the contact form $\theta(\xi) = 1$ singles out a
preferred direction along the vertical distribution.  Left-handed and
right-handed chiralities correspond to the two orientations of $\xi$ relative
to the complex structure of $M$.  For fermions
(Definition~\ref{def:fermion}), the vertical component $\alpha_V$ transforms
with a definite sign under reversal of $\xi$, so chirality is encoded in the
$\pm$-eigenspaces of the Reeb involution $\xi \mapsto -\xi$.  Bosons
(Definition~\ref{def:boson}), being purely horizontal, are chirality-blind ---
consistent with the observed parity invariance of the electromagnetic and
strong interactions.
\end{remark}

% ── Light cone and affine charts ─────────────────────────────────────────────

\begin{remark}[Light cone of an affine chart]
\label{rem:light-cone-affine}
Consider an open subset $U \subset M_0$ representing a portion of the
descended universe from the big bang to heat death.  An \emph{affine chart}
on $U$ is a local coordinate patch in which the metric on the base manifold
$M_0$ defines a causal structure.  The \emph{light cone} of the affine chart
is the boundary of the causal future in $U$: it is determined by the metric
on $M_0$ and has the property that no singularity of the ambient space is
visible from within the chart.  In this way, the affine chart provides a
non-singular window onto a portion of the Descartes universe.
\end{remark}

% ── Trefoil subgroups and Bring's curve ──────────────────────────────────────

\begin{remark}[Chirality and trefoil subgroups]
\label{rem:chirality-trefoil}
The chirality of the observable universe --- manifest in the three colour
charges of quantum chromodynamics --- is encoded in the symmetry group of the
descended surface $M_0$.  Proposition~\ref{prop:genus-maximal-descent}
identifies $M_0$ as a genus-$4$ Riemann surface.  We require a genus-$4$
surface whose automorphism group contains three subgroups isomorphic to $S_3$
--- quotients of the trefoil knot group
$\pi_1(S^3 \setminus K) \cong \langle x, y \mid x^2 = y^3 \rangle$ ---
corresponding to the three colour charges (red, green, blue).  These $S_3$
subgroups generate a subgroup of the monodromy group, and the remaining
generators of the full automorphism group act by conjugation, permuting the
three $S_3$ subgroups among themselves.
\end{remark}

\begin{remark}[Bring's curve]
\label{rem:brings-curve}
Bring's curve is the unique Riemann surface of genus $4$ with the maximal
automorphism group $S_5$ (order~$120$) for its genus.  It is defined by
the intersection of three symmetric hypersurfaces in $\mathbb{P}^4$:
\[
  \sum_{i=1}^{5} z_i = 0,
  \qquad
  \sum_{i=1}^{5} z_i^2 = 0,
  \qquad
  \sum_{i=1}^{5} z_i^3 = 0.
\]
The symmetric group $S_5$ acts by permuting the five homogeneous coordinates,
and this action realises the monodromy structure described in
Remark~\ref{rem:chirality-trefoil}: the group $S_5$ contains three copies of
$S_3$ (quotients of the trefoil knot group) corresponding to the three
colour charges, and elements outside these subgroups act by conjugation,
permuting the three copies among themselves.
\end{remark}

% ── Colour-sector coordinates on the ℂP¹ exceptional fibre ──────────────────

\begin{definition}[Exceptional divisor at $p$]
\label{def:exceptional-divisor}
Let $(M, G, E)$ be a Descartes universe and let $p \in M$ be a brane point
with outward normal $N(p) = p/\lvert p\rvert$ and positively-oriented
tangent frame $(T_1, T_2)$.  Set $R = R_{\mathbb{P}^1(\mathbb{C})}$ and
\[
  O(p) \;=\; p + R\,N(p).
\]
The \emph{exceptional divisor} at $p$ is the $\mathbb{P}^1(\mathbb{C})$ fibre
\[
  E_p \;=\; \pi^{-1}(p)
  \;=\;
  \bigl\{\,
    O(p) + R\bigl(\sin\phi\,(\cos\psi\,T_1 + \sin\psi\,T_2) - \cos\phi\,N\bigr)
    \;:\;
    \phi \in [0,\pi],\; \psi \in [0,2\pi)
  \,\bigr\},
\]
where $\pi\colon \widetilde{M} \to M$ is the blow-up map that replaces $p$
with $E_p \cong \mathbb{P}^1(\mathbb{C})$.
\end{definition}

\begin{definition}[Chiral sector coordinates $\mathfrak{r},\mathfrak{g},\mathfrak{b}$]
\label{def:chiral-sector-coords}
With notation as in Definition~\ref{def:exceptional-divisor}, fix the
substrate parameter $q$ (Definition~\ref{def:substrate-parameter}) and set
$\phi \in [0,\pi]$, $\psi \in [0, 2\pi/3]$.  Define the
\emph{red}, \emph{green}, and \emph{blue} sector embeddings by
\begin{align*}
  \mathfrak{r}(\phi,\psi)
  &\;=\; O + R\bigl(\sin\phi\,\bigl(\cos(\psi + q\phi)\,T_1
         + \sin(\psi + q\phi)\,T_2\bigr) - \cos\phi\,N\bigr),\\
  \mathfrak{g}(\phi,\psi)
  &\;=\; \mathfrak{r}\!\bigl(\phi,\,\psi + \tfrac{2\pi}{3}\bigr),\\
  \mathfrak{b}(\phi,\psi)
  &\;=\; \mathfrak{r}\!\bigl(\phi,\,\psi + \tfrac{4\pi}{3}\bigr).
\end{align*}
The \emph{chiral twist} $q\phi$ rotates each sector by $q$ full turns as
the latitude $\phi$ runs from the south pole~($0$) to the north pole~($\pi$),
encoding the chirality of the substrate parameter.

The three sectors tile the exceptional divisor exactly:
\[
  E_p \;=\; \mathfrak{r}(p) \;\cup\; \mathfrak{g}(p) \;\cup\; \mathfrak{b}(p),
\]
where adjacent sectors share only their boundary meridians
(at $\psi = 0$ and $\psi = 2\pi/3$ in each sector's local coordinate),
and the union covers $E_p$ without overlap.  The three sectors correspond to
the three copies of $S_3 \subset S_5$ in
Remark~\ref{rem:chirality-trefoil}: $\mathfrak{r}$ is the first, $\mathfrak{g}$ the
second, and $\mathfrak{b}$ the third.
\end{definition}

% ── Categorical braid structure of the Descartes universe ───────────────────

\begin{corollary}[Braid-group structure of a Descartes universe]
\label{cor:descartes-braid-structure}
Let $(\mathcal{C}, A, \omega)$ be a Descartes universe
(Definition~\ref{def:descartes-condition}).  Then:
\begin{enumerate}
  \item The braiding of $\mathcal{C} = \mathrm{Rep}(\mathrm{SU}_q(2))$
        induces a faithful representation
        $\rho_3\colon B_3 \to \mathrm{Aut}_{\mathcal{C}}(V^{\otimes 3})$
        on the fundamental representation~$V$
        (Proposition~\ref{prop:braid-action-tensor}).
  \item The image of $\rho_3$ in $\mathrm{SU}(2)$ (via the fibre
        functor~$\omega$) is the binary dihedral group
        $\mathrm{Dic}_3$ of order~$12$
        (Definition~\ref{def:su2-generators}).
  \item The composite $B_3 \xrightarrow{\;\rho_3\;}
        \mathrm{Dic}_3 \twoheadrightarrow S_3$ recovers the three
        copies of $S_3$ inside $S_5 = \mathrm{Aut}(\Sigma_4)$
        (Remark~\ref{rem:chirality-trefoil}), one for each
        colour sector.
  \item This structure is defined over any base ring~$R$ with
        $30 \in R^{\times}$; in particular it applies to the
        $\pp$-adic realisation
        (Example~\ref{ex:p-adic-realisation}), where knot complements
        and skein relations are not available.
\end{enumerate}
\end{corollary}

\begin{proof}
Item~(1) follows from Proposition~\ref{prop:braid-action-tensor}
applied to the fundamental representation $V$ of
$\mathrm{SU}_q(2)$, which is $2$-dimensional.  The braiding
$c_{V,V}$ is computed by the $R$-matrix of
$\mathcal{U}_q(\mathfrak{su}(2))$
(Remark~\ref{rem:r-matrix-braiding}), and faithfulness follows from the
fact that the eigenvalues of $c_{V,V}$ are $q$ and $-q^{-1}$, which
are distinct for $q \neq \pm 1$.

For item~(2), the fibre functor $\omega$ sends the braiding to the
matrix $\check{R} = P \circ \mathcal{R}$ acting on
$\mathbb{C}^2 \otimes \mathbb{C}^2$.  The generators $\sigma_1, \sigma_2$
map to the $\mathrm{SU}(2)$ elements $\alpha, \beta$
(Definition~\ref{def:su2-generators}) generating $\mathrm{Dic}_3$, as
verified in Example~\ref{ex:trefoil-operator-algebra}.

Item~(3): the projection $\mathrm{Dic}_3 \to S_3$ sends
$\alpha \mapsto (1\,2\,3)$ and $\beta \mapsto (1\,2)$.  The three
copies of $S_3$ inside $S_5$ are the stabilisers of the three pairs
of chiral coordinates
(Definition~\ref{def:chiral-sector-coords}), and conjugation by
elements of $S_5 \setminus S_3$ permutes these copies.

Item~(4): the braided monoidal structure of
$\mathrm{Rep}(\mathrm{SU}_q(2))$ is defined over~$R[q, q^{-1}]$
with $30 \in R^{\times}$
(Proposition~\ref{prop:brings-is-qou}(1)), and
Proposition~\ref{prop:braid-action-tensor} requires no topological
input beyond the hexagon axioms.
\end{proof}

\begin{remark}[Skein relations as an archimedean shadow]
\label{rem:skein-archimedean-shadow}
In the archimedean realisation, the Markov trace on the braid-group
representation $\rho_3$ yields the Jones polynomial, and the skein
relation of Chapter~\ref{chap:experimental} is the trace of the
characteristic equation of~$c_{V,V}$.  Over $\mathbb{Q}_\pp$, the
same characteristic equation holds (the eigenvalues $q$ and $-q^{-1}$
are elements of the base ring), but there is no ambient $3$-manifold in
which to interpret it as a knot invariant.  The categorical formulation
of Corollary~\ref{cor:descartes-braid-structure} therefore captures
the algebraic content of braiding without presupposing an archimedean
topology.
\end{remark}

% ── The Descartes theorem ────────────────────────────────────────────────────

\begin{theorem}[Descartes theorem]
\label{thm:descartes}
Let $\mathcal{C}_{-6}$ denote the category of quantum observable universes
over~$R$ satisfying the Descartes condition $\chi(A) = -6$
(Definition~\ref{def:descartes-condition}), with morphisms given by QOU maps
that preserve the Euler characteristic.  Then Bring's surface $\Sigma_4$
(Definition~\ref{def:brings-surface},
Proposition~\ref{prop:brings-is-qou}) is a terminal object in
$\mathcal{C}_{-6}$: for every Descartes universe
$(\mathcal{C}, A, \omega)$ in $\mathcal{C}_{-6}$, there exists a unique
morphism to the Bring's QOU.
\end{theorem}

\begin{proof}
\emph{Existence.}\;
By Theorem~\ref{thm:qu-to-qou} and
Proposition~\ref{prop:genus-maximal-descent}, the base curve is a smooth
genus-$4$ curve $\Sigma$ with $\chi(\Sigma) = -6$.
The Hurwitz bound gives $\lvert\mathrm{Aut}(\Sigma)\rvert \leq 252$;
among genus-$4$ curves, Bring's surface achieves the maximum
$\lvert S_5\rvert = 120$ (Edge, 1978), and the $S_5$-equivariant structure
uniquely determines $\Sigma_4 \subset \mathbb{P}^4_R$.
The descent data (branch locus with $6$ points, $S_5$-symmetric
ramification) defines a map $\Sigma \to \Sigma_4$, extending to a QOU
morphism by functoriality of the Haar state on
$\mathcal{O}_q(\mathrm{SU}(2))$.

\emph{Uniqueness.}\;
By Tannaka--Krein rigidity, any two fibre-preserving monoidal functors
between Tannakian categories are naturally isomorphic.
\end{proof}

\begin{remark}[Universality for a given Euler characteristic]
\label{rem:universality-general}
The Descartes theorem (Theorem~\ref{thm:descartes}) is expected to be a
special case of a more general result: for any admissible Euler
characteristic $\chi$, the category $\mathcal{C}_{\chi}$ of $\mathbb{C}$
quantum observable universes with $\chi(E) = \chi$ should admit a terminal
object determined by the unique maximally symmetric Riemann surface of the
corresponding genus.  The Descartes case $\chi = -6$ (genus~$4$) is
distinguished by Bring's curve, but analogous universal objects are expected
at other values of~$\chi$.
\end{remark}

\begin{example}[Bring's surface and Yekutieli obstructions]
\label{ex:descartes-yekutieli}
Bring's surface satisfies Theorem~\ref{thm:qu-to-qou} with
$\chi(A) = -6$.  The non-vanishing Yekutieli classes $\kappa_n \neq 0$
(Theorem~\ref{thm:qu-to-qou}(3)) obstruct smooth lifting at every level,
creating twisted sectors that control inertial mass
(Remark~\ref{rem:mass-geometric-drag}) and the three lepton generations
--- corresponding to the three pairs of branch points under
$S_3 \subset S_5$ (Remark~\ref{rem:chirality-trefoil}).
\end{example}

\begin{remark}[Observations on the Bring surface]
\label{rem:observation-brings}
The general notion of an observation
(Definition~\ref{def:observation}) specialises on the Bring surface as
follows: the observation morphism $\mathcal{O} \colon A \to A$ acts
fibrewise on the state bundle $E \to M$ and perturbs the substrate
parameter $q(p)$ (Definition~\ref{def:substrate-parameter}), shifting
the non-commutative geometry of the quantum Bring surface
(Section~\ref{sec:quantum-brings}).  The two-observer parallelogram
(Remark~\ref{rem:two-observer-parallelogram}) applies directly: two
independent observations of a fermion on the Bring surface produce a
discrepancy controlled by the local Planck constant $\hbar_q(p)$.
\end{remark}

